\documentclass[10pt,a4paper]{amsart}
\usepackage[margin=19mm]{geometry}
\usepackage{amsmath,amssymb,mathtools}
\usepackage[scr=rsfs,cal=euler]{mathalfa}
\usepackage{xfrac}
\usepackage[unicode,hidelinks,bookmarks=false]{hyperref}
\newcommand{\ie}{i.e.\ }
\newcommand{\cf}{cf.\ }
\newcommand{\resp}{respectively\ }
\newcommand{\Subsection}{\subsection}
\providecommand{\nobreakdash}{\nobreak\hbox{-}}
\setlength{\emergencystretch}{2em}
\multlinegap0pt
\usepackage{bm}
\usepackage{bbm}
\usepackage{dsfont}
\usepackage{xspace}
\usepackage{cancel}
\usepackage{mathtools}
\usepackage[shortlabels]{enumitem}
\usepackage{amsthm}
\makeatletter
\@ifundefined{theorem}{\newtheorem{theorem}{Theorem}}{}
\@ifundefined{lemma}{\newtheorem{lemma}[theorem]{Lemma}}{}
\makeatother
\renewcommand{\phi}{\varphi} 
\title[Outrunning the Omega Clock: A Singular Control Problem for D]{Outrunning the Omega Clock: A Singular Control Problem for Dividend Optimisation with Ruin and Time-in-Distress Default}
\author{Andi Bodnariu, Nils Engler, Neofytos Rodosthenous}
\date{Source: arXiv:2601.21705v1 (CC BY 4.0)}
\begin{document}
\maketitle

\noindent\emph{Source and licence.} Outrunning the Omega Clock: A Singular Control Problem for Dividend Optimisation with Ruin and Time-in-Distress Default. Version arXiv:2601.21705v1 (CC BY 4.0), distributed under the Creative Commons Attribution 4.0 International licence, as recorded in the arXiv licence metadata for this exact version and shown on the abstract page https://arxiv.org/abs/2601.21705v1. Full rights notice: licenses/arxiv-2601.21705-CC-BY-4.0.txt.\par\smallskip

\noindent\textit{Editorial scope.} Counted unit: the Lemma whose source label is \texttt{lemma:yl} in the article (source0.tex lines 1440--1463, source label \texttt{lemma:yl}), delivered together with the article's own complete original proof of it (source0.tex lines 1465--1660, one \texttt{proof} environment of 196 lines). No mathematics is written, repaired or re-derived: statement and proof are the article's own text line for line. Required context retained uncounted: value\_classical (lines 524--533); gamma\_def (lines 535--537); lemma\_useful\_results\_gamma (lines 539--544); lemma\_useful\_results\_Vb (lines 546--550); lemma\_useful\_results (lines 554--556); yupper\_def (lines 622--630); FBP\_lowya (lines 696--700); FBP\_lowyd (lines 696--700); def:delta (lines 704--713); sol:FBP1 (lines 717--746); constants\_simple (lines 728--737); def\_deltay (lines 739--744); eq:C4divC3 (lines 1365--1370). Every definition, display and label of those lines is retained unchanged. Omitted from this package and not redistributed: 1--523, 534--534, 538--538, 545--545, 551--553, 557--621, 631--695, 701--703, 714--716, 738--738, 745--1364, 1371--1439, 1464--1464, 1661--2172. Nothing inside a retained excerpt is omitted or altered: every retained body is delivered exactly as the article's lines are written, so no omission receipt of any kind accompanies this package. Citations: the retained excerpts carry no citation command at all, so no bibliography entry of the article is transcribed and no reference is redistributed. Reference adaptation: none was needed and none was made; every reference inside the delivered unit resolves inside it, so no unresolved reference or citation is produced. Printed slips kept literal: no printed word, symbol, hypothesis or number of the counted statement or of its proof was altered. Layout and class: the article's own class is \texttt{article} with packages including inputenc, geometry, graphicx, float, todonotes, booktabs, array, paralist, amsmath, amssymb, amsthm, xifthen, xcolor, ulem; the wrapper is the lane's pinned \texttt{amsart} editorial wrapper at 10pt on A4, which restates the theorem-like environments the retained text needs and adds the article's own macro definitions that the retained text needs (\texttt{\textbackslash phi}), with extra wrapper packages (bm, bbm, dsfont, xspace, cancel, mathtools, enumitem). The delivered numbering is the unit's own. Assets: no retained span contains a figure environment, a graphics-inclusion command, a PDF-page inclusion, a table or a TikZ picture; no image file is copied into this package, the package declares no asset dependency and no image file is redistributed with it. Licence: version arXiv:2601.21705v1 of this article is distributed under the Creative Commons Attribution 4.0 International licence, as recorded in the arXiv licence metadata for this exact version and shown on the abstract page https://arxiv.org/abs/2601.21705v1. A full rights notice is delivered at licenses/arxiv-2601.21705-CC-BY-4.0.txt. Characters: every retained excerpt is pure ASCII, so the delivered engine, XeLaTeX with its default Latin Modern text font, reports no missing character.
\begin{align} \label{value_classical}
V_\rho(x) = \begin{cases}
\dfrac{e^{\gamma_1(\rho) x} - e^{\gamma_2(\rho)x}}{\gamma_1(\rho) e^{\gamma_1(\rho) b^*_\rho} - \gamma_2(\rho) e^{\gamma_2(\rho)b^*_\rho}}, 
\quad &x\in[0,b^*_\rho),\\
V_\rho(b^*_\rho) + x - b^*_\rho, 
\quad &x\in [b^*_\rho,\infty),
\end{cases}
\quad \text{and} \quad 
b_\rho^* = \frac{\log\left(\gamma_2^2(\rho) / \gamma_1^2(\rho)\right)}{\gamma_1(\rho)-\gamma_2(\rho)} ,
\end{align}

\begin{align} \label{gamma_def}
\gamma_1(\rho)= \sqrt{\frac{\mu^2}{\sigma^4}+\frac{2\rho}{\sigma^2}}-\frac{\mu}{\sigma^2},\quad \gamma_2(\rho)= -\sqrt{\frac{\mu^2}{\sigma^4}+\frac{2\rho}{\sigma^2}}-\frac{\mu}{\sigma^2}
\end{align}

\begin{align} \label{lemma_useful_results_gamma} 
\gamma_1(\rho+q) + \gamma_2(\rho + q)
= \gamma_1(\rho) + \gamma_2(\rho)
= \tfrac{\mu}{\rho}  \gamma_1(\rho) \gamma_2(\rho), 
\quad q > 0,
\end{align}

\begin{align} \label{lemma_useful_results_Vb}
V_{\rho}(b^*_{\rho})
=\frac{\gamma_2^2(\rho)-\gamma_1^2(\rho)}{\gamma_2^2(\rho)\gamma_1(\rho)-\gamma_1^2(\rho)\gamma_2(\rho)}
= \frac{\gamma_1(\rho) + \gamma_2(\rho)}{\gamma_1(\rho) \gamma_2(\rho)} = \frac{\mu}{\rho}.
\end{align}

\begin{lemma} \label{lemma_useful_results}
The mapping $\rho \mapsto b^*_{\rho}$ is strictly decreasing on $(0,\infty)$. 
\end{lemma} 

\begin{align} \label{yupper_def}
\begin{cases}
0\leq 0, & x \in (0,b^*_{r+q}], \\
\mu - (r+q) (x - b^*_{r+q} + \frac{\mu}{r+q})  \leq 0, & x \in (b^*_{r+q},y), \\
\mu - r (x - b^*_{r+q} + \frac{\mu}{r+q})  \leq 0, & x \in [y,\infty),
\end{cases} 
\quad \Leftrightarrow \quad 
y \geq y_u := b^*_{r+q} + \frac{\mu}r - \frac{\mu}{r+q},
\end{align}

\begin{align}
&\tfrac12 \sigma^2 w''(x;y)+\mu w'(x;y) - (r+qI_{\{x< y\}})w(x;y) = 0, \quad x\in (0,y)\cup(y,b^*(y)) \label{FBP_lowya}\\
&w(x;y) = w(b^*(y);y) + x - b^*(y), \qquad \qquad \qquad \qquad \quad x\in [b^*(y),\infty) \label{FBP_lowyb}\\ &w(0+;y)=0, \label{FBP_lowyc}\\ 
&w(\cdot;y) \in \mathcal{C}^2 \big((0,y) \cup (y,\infty) \big) \cap  \mathcal{C}^1(0,\infty) \label{FBP_lowyd}.
\end{align}

\begin{align}
\label{def:delta}
\begin{split}
\delta(y) := &\psi_{r+q}(y) - \phi_{r+q}(y), \quad 
\eta(y;b) := \psi'_r(b)\phi_r(y) - \phi'_r(b) \psi_r(y), \quad \\ 
&\psi_\rho(y) := e^{\gamma_1(\rho) \,y}, \quad 
\phi_\rho(y) := e^{\gamma_2(\rho) \,y}, \quad 
y, b, \rho > 0 . 
\end{split}
\end{align}

\begin{lemma} \label{sol:FBP1}
For any $y>0$, the free-boundary problem \eqref{FBP_lowya}--\eqref{FBP_lowyd} admits a unique solution given by 
\begin{align} \label{w_sol_simple_def}
w(x;y)= 
\begin{cases}
K_1(y) \, e^{\gamma_1(r+q) \, x} + K_2(y) \, e^{\gamma_2(r+q) \, x}, \quad &x \in [0,y),\\
K_3(y) \, e^{\gamma_1(r) \, x} + K_4(y) \, e^{\gamma_2(r) \, x}, \quad &x \in [y,b^*(y)),\\
w(b^*(y);y) + x - b^*(y), \quad &x\in [b^*(y),\infty),
\end{cases}
\end{align}
with $\gamma_1(\cdot)$, $\gamma_2(\cdot)$ defined by \eqref{gamma_def}, $K_1(y), \ldots, K_4(y)$ defined by 
\begin{align}\label{constants_simple}
\begin{split}
K_1(y) = & - K_2(y) =  
\frac{\psi_r(y) \eta' (y; b^*(y)) - \psi'_r(y)\eta (y; b^*(y))}{\psi'_r(b^*(y))\bigl(\delta(y) \eta'(y; b^*(y)) - \delta'(y) \eta(y; b^*(y))\bigr)}, \\
K_3(y) = 
&\frac{\phi'_r(b^*(y)) \bigl( \delta \psi'_r - \delta' \psi_r \bigr)(y) + \delta(y) \eta'(y; b^*(y)) - \delta'(y) \eta(y; b^*(y))}{\psi'_r(b^*(y))\bigl(\delta(y) \eta'(y; b^*(y)) - \delta'(y) \eta(y; b^*(y))\bigr)}, \\ 
K_4(y) = 
&\frac{ \big(\psi_r \delta' -  \psi'_r \delta \big)(y)}{\delta(y) \eta'(y; b^*(y)) - \delta'(y) \eta(y; b^*(y))},
\end{split}
\end{align}
where $\delta$, $\eta$, $\psi_\rho$ and $\phi_\rho$ are defined by \eqref{def:delta}, and 
\begin{align} \label{def_deltay}
\begin{split}
b^*(y) &=  y + \Delta(y), \quad \text{with} \quad  
\Delta(y) := \frac{1}{\gamma_1(r)-\gamma_2(r)}\log\left(\frac{(\delta'(y)-\gamma_1(r)\delta(y)) \, \gamma_2^2(r)}{(\delta'(y)-\gamma_2(r)\delta(y)) \, \gamma_1^2(r)}\right),
\end{split}
\end{align}
such that $\Delta(y) > 0$ for all $0 < y \leq y_u$ and $y \mapsto b^*(y)$ is strictly increasing on $(0,\infty)$.
\end{lemma} 

\begin{align}\label{eq:C4divC3}
\hspace{-1mm}
-\frac{K_4(y)}{K_3(y)}
= \frac{\delta\psi'_r - \delta'\psi_r}{\delta \phi'_r - \delta' \phi_r}(y) 
= e^{(\gamma_1(r) - \gamma_2(r))y} \frac{e^{\gamma_1(r+q)y}\bigl(\gamma_1(r+q) - \gamma_1(r)\bigr) + e^{\gamma_2(r+q)y}\bigl(\gamma_1(r) - \gamma_2(r+q)\bigr)}{e^{\gamma_1(r+q)y}\bigl(\gamma_1(r+q) - \gamma_2(r)\bigr) + e^{\gamma_2(r+q)y}\bigl(\gamma_2(r) - \gamma_2(r + q)\bigr)},
\end{align}

\begin{lemma}\label{lemma:yl}
Recall $b^*_{r+q}$ from \eqref{value_classical}, $\delta$ from \eqref{def:delta}, $K_1, \ldots, K_4$ from \eqref{constants_simple}, and $b^*, \Delta$ from \eqref{def_deltay}. 
Then, we have:
\begin{enumerate}[\rm (i)]
\item 
$b^*_{r+q} < b^*_{r} < b^*(y)$ for all $y > 0$;
\item 
$b^*_{r+q}$ is the unique $y$-value that satisfies $\delta''(y)=0$, 
and 
$\delta''(y) > 0$ if and only if $y > b^*_{r+q}$;
\item 
$\frac{\delta'(x)}{\delta'(b^*_{r+q})}>1$ for all $x\in (b^*_{r+q},\infty)$ 
and 
$\frac{\delta(y_u)}{\delta'(y_u)}< \frac{\mu}{r} < \frac{\delta(y_u)}{\delta'(b^*_{r+q})}$;
\item $\Delta(y_u)>0$ and $-1 < \Delta'(y) < 0$ for all $y > 0$;
\item $K_3(y) > 0$ and $K_4(y) < 0$ for all $y > 0$;
\item $K_1(y)$ admits the form
\begin{align}\label{k1_form} 
K_1(y)= \left( \left(\frac{\gamma_2^2(r)}{\gamma_1^2(r)}\frac{\delta'(y) - \gamma_1(r)\delta(y)}{\delta'(y) - \gamma_2(r)\delta(y)}  \right)^{\frac{\gamma_1(r)}{\gamma_1(r) - \gamma_2(r)}}
\left(\frac{\gamma_1(r)}{-\gamma_2(r)}\delta'(y) + \gamma_1(r)\delta(y) \right)\right)^{-1}
\end{align}
satisfying $K_1(y) >0$ for all $y > 0$, $K_1'(y) < 0$ for $b^*_{r+q} < y \leq y_u$, as well as $K_1(y_u)\delta'(b^*_{r+q})<1$.
\end{enumerate}
\end{lemma}

\begin{proof}[Proof of Lemma \ref{lemma:yl}]
We prove each part separately. 

\vspace{1mm}
{\it Proof of part} (i).
The first inequality $b^*_{r+q} < b^*_r$ is proved in Lemma \ref{lemma_useful_results}, thus it suffices to prove that $b^*_r < b^*(y)$ for all $y > 0$. 
To that end, we use their expressions in \eqref{value_classical} and \eqref{def_deltay} to get  
\begin{align}\label{eq:bstar05}
\begin{split}    
b^*(y) &:= y +\Delta(y) 
= \; y + b^*_r + \frac{1}{d(r)} \log\Big( \frac{c_1 e^{\gamma_1(r+q)y} + c_2 e^{\gamma_2(r+q)y}}{c_2 e^{\gamma_1(r+q)y} + c_1 e^{\gamma_2(r+q)y}} \Big),
\end{split}    
\end{align}
for the positive constants $c_1, c_2$ and the increasing function $d(\cdot)$ defined by 
\begin{align}\label{cd_def}
\begin{split}
0 < c_1 := \gamma_1(r+q) - \gamma_1(r) = \gamma_2(r) - \gamma_2(r + q) &< \gamma_1(r) - \gamma_2(r+q) = \gamma_1(r+q) - \gamma_2(r) =: c_2,\\
0 < d(r) := \gamma_1(r) - \gamma_2(r) = c_2 - c_1 &< c_2 + c_1 = \gamma_1(r+q) - \gamma_2(r+q) =: d(r+q),
\end{split}
\end{align}
where the equalities in the first line follow from \eqref{lemma_useful_results_gamma}.
Then, we observe from \eqref{eq:bstar05}--\eqref{cd_def} that
\begin{align*}
\begin{split}
b^*_r < b^*(y)
\quad \Leftrightarrow \quad 
L(y) := e^{- d(r) y} \, \frac{c_2 e^{d(r+q)y} + c_1}{c_1 e^{d(r+q)y} + c_2} < 1 ,
\quad y>0 .
\end{split}    
\end{align*}
This holds true since $L(0)=1$ and $L'(y)<0$ for all $y>0$, which follows from \eqref{eq:bstar05} and 
\begin{align*}
\big(e^{d(r+q)y}c_1 + c_2 \big)^2 L'(y) 
&= - c_1 c_2 d(r) e^{-d(r)y} \big(e^{2d(r+q)y} + 1 \big) \\ 
&\quad + e^{(d(r+q)-d(r))y} \big(c_2^2 (d(r+q)-d(r)) -c_1^2 (d(r+q)+d(r)) \big) \\
&= -c_1 c_2 d(r) e^{-d(r)y} \big(e^{2d(r+q)y} - 2e^{d(r+q)y} + 1 \big) < 0.
\end{align*}
The latter positivity follows by observing that $\widetilde L(y):=e^{2d(r+q)y} - 2e^{d(r+q)y} + 1 > 0$ for all $y>0$, thanks to $\widetilde L(0)=0$ and $\widetilde L'(y) = 2d(r+q)\left(e^{2d(r+q)y} -e^{d(r+q)y}\right) > 0$ for $y > 0$.

\vspace{1mm}
{\it Proof of part} (ii).
Given that 
$\delta''(y) = \gamma_1^2(r+q)e^{\gamma_1(r+q)y} - \gamma_2^2(r+q)e^{\gamma_2(r+q)y}$, we can directly verify that $y = b^*_{r+q}$ is the only solution to $\delta''(y)=0$ and that $\delta''(y)>0$ for all $y > b^*_{r+q}$. 

\vspace{1mm}
{\it Proof of part} (iii).
First note that the following equalities hold true (cf.~\eqref{value_classical}, \eqref{yupper_def}) 
\begin{align*} 
V'_{r+q}(x)=1 \quad \text{for} \quad x \geq b^*_{r+q} 
\quad \text{with} \quad 
V_{r+q}(b^*_{r+q}) = \tfrac{\delta(b^*_{r+q})}{\delta'(b^*_{r+q})} 
\quad \text{and} \quad 
V_{r+q}(y_u)=\frac{\mu}{r}. 
\end{align*}
Combining the above with 
the fact that $\delta''(x)>0$, for all $x > b^*_{r+q}$, thanks to Lemma \ref{lemma:yl}.(ii), and the fundamental theorem of calculus, we obtain the desired results via 
\begin{align*}
\frac{\mu}{r} - \frac{\delta(y_u)}{\delta'(y_u)} 
&= V_{r+q}(y_u) - V_{r+q}(b^*_{r+q}) - \frac{\delta(y_u)}{\delta'(y_u)} + \frac{\delta(b^*_{r+q})}{\delta'(b^*_{r+q})} 
= \int_{b^*_{r+q}}^{y_u} \Big( V_{r+q}'(x) - 1 + \frac{\delta(x)\delta''(x)}{(\delta'(x))^2} \Big) dx > 0, \notag\\
\frac{\delta'(x)}{\delta'(b^*_{r+q})} - 1 &= \int_{b^*_{r+q}}^{x} \frac{\delta''(z)}{\delta'(b^*_{r+q})} dz > 0, \quad x>b^*_{r+q}, \\
\frac{\delta(y_u)}{\delta'(b^*_{r+q})} - \frac{\mu}{r} 
&= \frac{\delta(y_u)}{\delta'(b^*_{r+q})} - \frac{\delta(b^*_{r+q})}{\delta'(b^*_{r+q})} + V_{r+q}(b^*_{r+q}) - V_{r+q}(y_u) = \int_{b^*_{r+q}}^{y_u} \Big(\frac{\delta'(x)}{\delta'(b^*_{r+q})} - 1 \Big) dx > 0.
\end{align*}

\vspace{1mm}
{\it Proof of part} (iv).
Firstly, we recall from Lemma \ref{lemma:yl}.(iii) and the expression of $\gamma_1$ from \eqref{gamma_def} that
\begin{align} \label{ineq1}
\frac{\delta(y_u)}{\delta'(y_u)}<\frac{\mu}{r} < \frac1{\gamma_1(r)} 
\quad \Leftrightarrow \quad 
\delta'(y_u)-\gamma_1(r)\delta(y_u) > 0.
\end{align}
Combining this with the definition \eqref{def_deltay} of $\Delta(y)$ and the positivity of $\delta(y), \delta'(y)$ from \eqref{def:delta} for all $y>0$, together with the fact that $\gamma_2(r) < 0 < \gamma_1(r)$, we get that
\begin{align*}
&\Delta(y_u) 
= \frac{\log\left( \frac{(\delta'(y_u)-\gamma_1(r)\delta(y_u)) \gamma_2^2(r)}{(\delta'(y_u)-\gamma_2(r)\delta(y_u)) \gamma_1^2(r)} \right)}{\gamma_1(r)-\gamma_2(r)}  > 0 
\quad \Leftrightarrow \quad 
(\delta'(y_u)-\gamma_1(r)\delta(y_u))\gamma_2^2(r)>(\delta'(y_u)-\gamma_2(r)\delta(y_u))\gamma_1^2(r).
\end{align*}
To see how the latter inequality holds true, divide both sides by 
$-\gamma_1(r) \gamma_2(r) \delta'(y_u) > 0$, and observe that it is equivalent to
\begin{align*}
\frac{\gamma_2^2(r)-\gamma_1^2(r)}{-\gamma_1(r)\gamma_2(r)} - \frac{\delta(y_u)}{\delta'(y_u)}(\gamma_1(r)-\gamma_2(r)) >0,
\end{align*}
which holds true thanks to $\frac{\delta(y_u)}{\delta'(y_u)}<\frac{\mu}{r}$ from \eqref{ineq1} (see also Lemma~\ref{lemma:yl}.(iii)) and \eqref{lemma_useful_results_Vb}. 

For the monotonicity of $\Delta(\cdot)$ defined by \eqref{def_deltay}, 
which takes the form as in \eqref{eq:bstar05}
\begin{align*} 
\begin{split}
\Delta(y) =b^*_r+ \frac{1}{d(r)}\log \Big(\frac{u(y)}{v(y)} \Big),
\quad \text{where} \quad 
u(y):=c_1e^{\gamma_1(r+q)y} + c_2e^{\gamma_2(r+q)y},  
\quad 
v(y):=c_2e^{\gamma_1(r+q)y} + c_1e^{\gamma_2(r+q)y},
\end{split}
\end{align*}
for $0 < c_1 <  c_2 < \infty$, $0 < d(r) < d(r+q) < \infty$ given by \eqref{cd_def}, 
we calculate 
\begin{align} \label{eq:DeltaPrime1}
\Delta'(y) 
= \frac{u'(y) v(y) - u(y) v'(y)}{d(r)u(y)v(y)} < 0.
\end{align}
The latter inequality follows from the straightforward positivity of all three factors in the denominator, and the negativity of the numerator which results from 
\begin{align}\label{eq:DeltaPrime2}
\begin{split}
u'(y) v(y) - u(y) v'(y) 
&= \big(\gamma_1(r+q) c_1 e^{\gamma_1(r+q)y} + \gamma_2(r+q) c_2 e^{\gamma_2(r+q)y}\big) \big(c_2 e^{\gamma_1(r+q)y} + c_1 e^{\gamma_2(r+q)y} \big) \\ 
&\quad - \big(c_1 e^{\gamma_1(r+q)y} +  c_2 e^{\gamma_2(r+q)y} \big) \big(\gamma_1(r+q) c_2 e^{\gamma_1(r+q)y} + \gamma_2(r+q) c_1 e^{\gamma_2(r+q)y} \big) \\
&= e^{(\gamma_1(r+q) +  \gamma_2(r+q))y} \, d(r+q) \left(c_1^2 - c_2^2 \right) < 0 ,
\end{split}
\end{align}
thanks to $0 < c_1 < c_2$. 
To complete the monotonicity properties, it remains to show that $\Delta'(y) > -1$ for all $y > 0$. 
To that end, using \eqref{cd_def}, \eqref{eq:DeltaPrime1} and \eqref{eq:DeltaPrime2}, we get 
\begin{align*}
\Delta'(0) 
= \frac{d(r+q)(c_1^2 - c_2^2)}{d(r) (c_1 + c_2)^2} = -1,
\quad \text{since} \quad 
c_1^2 - c_2^2 = -d(r)d(r+q)
\quad \text{and} \quad 
(c_1+c_2)^2 = d(r+q)^2 .
\end{align*}
Combining this with $\Delta \in C^2([0, \infty))$ and \eqref{eq:DeltaPrime1}--\eqref{eq:DeltaPrime2}, we get
\begin{align*}
\Delta''(y) 
&= - d(r+q)^2 \frac{d}{dy} \Big(\frac{e^{(\gamma_1(r+q) + \gamma_2(r+q))y}}{u(y)v(y)} \Big) = d(r+q)^3 c_1 c_2 \frac{\big(e^{2\gamma_1(r+q)y} - e^{2\gamma_2(r+q)y} \big) e^{(\gamma_1(r+q)+\gamma_2(r+q))y}}{u^2(y) v^2(y)} > 0,
\end{align*}
which implies that $\Delta'(\cdot)$ is strictly increasing, thus $\Delta'(y) > -1$ for all $y > 0$.

\vspace{1mm}
{\it Proof of part} (v). 
Recall the definitions \eqref{constants_simple} of $K_3$ and $K_4$ and \eqref{cd_def} of the constants $0 < c_1 <  c_2 < \infty$, as well as the fact that $\gamma_2(r) < 0 < \gamma_1(r)$. 
Substituting the expression \eqref{def_deltay} of $b^*$ in \eqref{constants_simple} then gives
\begin{align*}
K_4(y) &= -\frac{e^{-\gamma_2(r)y} \big(e^{\gamma_1(r+q)y} c_1 + e^{\gamma_2(r+q)y} c_2 \big)
    }
    {
    e^{\gamma_1(r+q)y} \left( 
    e^{\gamma_1(r)\Delta(y)} c_2\gamma_1(r) - e^{\gamma_2(r)\Delta(y)} c_1\gamma_2(r)
    \right) 
    + e^{\gamma_2(r+q)y} \left(
    e^{\gamma_1(r)\Delta(y)} c_1\gamma_1(r) - e^{\gamma_2(r)\Delta(y)} c_2\gamma_2(r) \right)
    },
\end{align*}
which is negative for all $y > 0$. 
Combining this with the positivity of the fraction in \eqref{eq:C4divC3}, we can conclude that $K_3(y)$ has the opposite sign to $K_4(y)$, hence $K_3(y) > 0$ for $y > 0$.

\vspace{1mm}
{\it Proof of part} (vi). 
Using the definition of $K_1$ from \eqref{constants_simple}, we observe after long calculations that 
\begin{align}
K_1(y) &= \frac{\gamma_1(r)-\gamma_2(r)}{G(y)}, \quad \text{where} \quad 
\label{k_1_repg}\\ 
G(y)&:= \left( \gamma_1(r) e^{\gamma_1(r)\Delta(y)}  - \gamma_2(r) e^{\gamma_2(r)\Delta(y)}\right)\delta'(y) - \gamma_1(r)\gamma_2(r)\left(e^{\gamma_1(r) \Delta(y)} - e^{\gamma_2(r)\Delta(y)} \right) \delta(y) 
\label{eq:G_y}.
\end{align}
Using \eqref{def_deltay} then gives 
\begin{align}
G(y)
&=e^{\gamma_1(r)\Delta(y)}\left(\left( \gamma_1(r)  - \gamma_2(r) e^{(\gamma_2(r)-\gamma_1(r))\Delta(y)}\right)\delta'(y) - \gamma_1(r)\gamma_2(r)\left(1 - e^{(\gamma_2(r)-\gamma_1(r))\Delta(y)} \right) \delta(y)\right)\notag\\
&=e^{\gamma_1(r)\Delta(y)} \left(-\gamma_2(r)e^{(\gamma_2(r)-\gamma_1(r))\Delta(y)}\left(\delta'(y)-\gamma_1(r)\delta(y)\right)+\gamma_1(r)\delta'(y)-\gamma_1(r)\gamma_2(r)\delta(y)\right)\notag\\
&=e^{\gamma_1(r)\Delta(y)}\Big(-\frac{\gamma_1^2(r)}{\gamma_2(r)}(\delta'(y)-\gamma_2(r)\delta(y))+\gamma_1(r)\delta'(y)-\gamma_1(r)\gamma_2(r)\delta(y)\Big)\notag\\
&=e^{\gamma_1(r)\Delta(y)}(\gamma_1(r)-\gamma_2(r))\Big(\gamma_1(r)\delta(y)-\frac{\gamma_1(r)}{\gamma_2(r)}\delta'(y)\Big) > 0, \quad y > 0,
\label{g_rep1} 
\end{align}
whose positivity follows from \eqref{g_rep1} thanks to the fact that $\gamma_2(r) < 0 < \gamma_1(r)$ and $\delta(y), \delta'(y) > 0$ for $y>0$ by definition \eqref{def:delta}.
This provides the desired form of $K_1$ in \eqref{k1_form} upon substituting the expression of $\Delta$ from \eqref{def_deltay} into \eqref{g_rep1}.
Hence, the positivity of $K_1(y)$ for all $y>0$ follows from \eqref{k_1_repg} and \eqref{g_rep1}.

Regarding the monotonicity of $K_1$, we notice that $K_1'(y) = - (\gamma_1(r)-\gamma_2(r)) \frac{G'(y)}{G^2(y)}$, whose sign is the opposite of that of $G'(y)$. Taking the derivative of $G$ in \eqref{eq:G_y} gives
\begin{align} \label{gprime_eq}
G'(y)
&= \Delta'(y) \left[ \gamma_1^2(r)e^{\gamma_1(r)\Delta(y)}\left( \delta'(y) - \gamma_2(r) \delta(y) \right)
- \gamma_2^2(r)e^{\gamma_2(r)\Delta(y)}\left( \delta'(y) - \gamma_1(r)\delta(y) \right)
\right] \notag\\
&\quad -\gamma_1(r)\gamma_2(r) \delta'(y)\left(e^{\gamma_1(r)\Delta(y) } - e^{\gamma_2(r)\Delta(y)} \right)
+ \delta''(y)\left(\gamma_1(r)e^{\gamma_1(r)\Delta(y) } - \gamma_2(r)e^{\gamma_2(r)\Delta(y)} \right) \notag\\
&= -\gamma_1(r)\gamma_2(r) \delta'(y)\left(e^{\gamma_1(r)\Delta(y) } - e^{\gamma_2(r)\Delta(y)} \right)
+ \delta''(y)\left(\gamma_1(r)e^{\gamma_1(r)\Delta(y) } - \gamma_2(r)e^{\gamma_2(r)\Delta(y)} \right),
\end{align}
where the second equality follows from the definition \eqref{def_deltay} of $\Delta(y)$. 
To see the positivity of $G'(y)$ for all $y \in (b^*_{r+q}, y_u]$, observe in \eqref{gprime_eq} that $\Delta(y) > 0$ thanks to Lemma \ref{sol:FBP1}, $\delta'(y) > 0$ thanks to \eqref{def:delta}, and $\delta''(y) > 0$ by Lemma \ref{lemma:yl}.(ii). 

Finally, using \eqref{k_1_repg} together with the positivity and the explicit expression of $G$ from \eqref{g_rep1}, we notice that $K_1(y_u)\delta'(b^*_{r+q})<1$ is equivalent to 
\begin{align*}
\delta'(b^*_{r+q}) < \frac{G(y_u)}{\gamma_1(r)-\gamma_2(r)} = e^{\gamma_1(r)\Delta(y_u)} \left(-\frac{\gamma_1(r)}{\gamma_2(r)}\delta'(y_u)+\gamma_1(r)\delta(y_u)\right). 
\end{align*}
To prove this, given that $\gamma_1(r)>0$ and $\Delta(y_u) > 0$ by Lemma \ref{lemma:yl}.(iv), it is sufficient to show that
\begin{align*}
- \frac{\gamma_1(r) \delta'(y_u)}{\gamma_2(r) \delta'(b^*_{r+q})} + \gamma_1(r) \frac{\delta(y_u)}{\delta'(b^*_{r+q})} - 1 > 0.
\end{align*}
This follows from $\gamma_2(r)<0< \gamma_1(r)$ in \eqref{gamma_def}, Lemma \ref{lemma:yl}.(iii), \eqref{yupper_def} and \eqref{lemma_useful_results_gamma}, 
and completes the proof.
\end{proof}

\end{document}
